\section{Introduction}

A tunnel endpoint that terminates one security association and originates
another pays two kinds of cost for every packet it forwards. One kind does not
depend on how large the packet is: the system calls that move it across the
kernel boundary, the header parse, the anti-replay check of ESP and its
descendants \cite{rfc4303}, the nonce derivation, and the cipher's own
per-message setup and teardown. The other kind is proportional to the payload,
and covers the bulk encryption and decryption together with the memory traffic
that feeds them.

Writing the first as $a$ cycles and the second as $b$ cycles per payload byte
gives
%
\begin{equation}
C(s) = a + b\,s ,
\qquad
s^{*} = \frac{a}{b} ,
\label{eq:model}
\end{equation}
%
where $s^{*}$ is the payload at which the two halves are equal. Below $s^{*}$ a
packet costs mostly what it costs to be a packet at all; above it, mostly what
it costs to carry that many bytes. The two sides call for opposite work. Below
$s^{*}$, the useful optimisations are batching and the removal of system calls.
Above it, they are a faster cipher or fewer passes over the payload.

The ratio is more portable than either term. Both $a$ and $b$ are counts of
cycles, and both scale with the clock, so $s^{*}$ survives clock scaling and can
be compared across machines whose absolute numbers cannot. That portability is
why this study reports it.

\noindent\textbf{What is measured here.}
The data plane is a decrypt-then-re-encrypt forwarding element, so every payload
byte crosses the cipher twice. Crossing twice makes $b$ large enough to separate
cleanly from $a$, and the report states the factor of two rather than folding it
away. Four factors are varied against it: the cipher, in four levels; the
presence of the hardware crypto instructions; the transport, in four levels
whose system-call rates differ by construction; and the payload size, over a
grid dense at the small end where $a$ dominates.

\noindent\textbf{Two implementations of each cipher.}
For AES-256-GCM \cite{nistgcm} and ChaCha20-Poly1305 \cite{rfc8439} we measure
OpenSSL's implementation and one written for this study. Ours uses AES-NI with a
PCLMULQDQ GHASH \cite{gueron2010} and a stitched single-pass
counter-mode-and-hash inner loop \cite{krasnov2017} on x86-64, ARMv8 AES with
PMULL on \texttt{aarch64}, and ChaCha20 vectorised with AVX2 or NEON. Both agree
with OpenSSL byte for byte at every length in the sweep, which the self-test
checks before any timing run. A second implementation is what separates ``this
cipher costs $x$'' from ``this library's version of this cipher costs $x$'', and
the gap between the two is reported rather than hidden.

\noindent\textbf{Contributions.}
\begin{itemize}
  \item A measured cost model for an encrypted UDP data plane, with the
        crossover payload $s^{*}$ and its confidence interval, for four cipher
        implementations with hardware crypto present and masked.
  \item A direct measurement of what the system-call boundary contributes to
        $a$, taken from a submission path that issues no system calls
        \cite{iouring} rather than extrapolated from a batching curve.
  \item The same fixed-versus-marginal decomposition applied to two further
        boundaries: the isolation layer around a virtual network function, and
        the bus between a CPU and a GPU.
  \item An artefact that refuses to produce a measurement on a host that cannot
        hold one, records why every unavailable experiment was unavailable, and
        generates every number in this text from the dataset.
\end{itemize}
